\documentclass[10pt,a4paper]{amsart}
\usepackage[margin=19mm]{geometry}
\usepackage{amsmath,amssymb,mathtools}
\usepackage[scr=rsfs,cal=euler]{mathalfa}
\usepackage{xfrac}
\usepackage[unicode,hidelinks,bookmarks=false]{hyperref}
\newcommand{\ie}{i.e.\ }
\newcommand{\cf}{cf.\ }
\newcommand{\resp}{respectively\ }
\newcommand{\Subsection}{\subsection}
\providecommand{\nobreakdash}{\nobreak\hbox{-}}
\setlength{\emergencystretch}{2em}
\multlinegap0pt
\usepackage{bm}
\usepackage{bbm}
\usepackage{dsfont}
\usepackage{xspace}
\usepackage{cancel}
\usepackage{mathtools}
\usepackage{amsthm}
\makeatletter
\@ifundefined{theorem}{\newtheorem{theorem}{Theorem}}{}
\@ifundefined{lemma}{\newtheorem{lemma}[theorem]{Lemma}}{}
\makeatother
\title[Adaptive Filtering via Canonical Systems with Time-Varying H]{Adaptive Filtering via Canonical Systems with Time-Varying Hamiltonians}
\author{Keshav Raj Acharya, Pitambar Acharya}
\date{Source: arXiv:2603.10096v1 (CC BY 4.0)}
\begin{document}
\maketitle

\noindent\emph{Source and licence.} Adaptive Filtering via Canonical Systems with Time-Varying Hamiltonians. Version arXiv:2603.10096v1 (CC BY 4.0), distributed under the Creative Commons Attribution 4.0 International licence, as recorded in the arXiv licence metadata for this exact version and shown on the abstract page https://arxiv.org/abs/2603.10096v1. Full rights notice: licenses/arxiv-2603.10096-CC-BY-4.0.txt.\par\smallskip

\noindent\textit{Editorial scope.} Counted unit: the Theorem whose source label is \texttt{None} in the article (source0.tex lines 327--333, source label \texttt{None}), delivered together with the article's own complete original proof of it (source0.tex lines 335--440, one \texttt{proof} environment of 106 lines). No mathematics is written, repaired or re-derived: statement and proof are the article's own text line for line. Required context retained uncounted: ceq (lines 262--264); lemma 1 (lines 289--299). Every definition, display and label of those lines is retained unchanged. Omitted from this package and not redistributed: 1--261, 265--288, 300--326, 334--334, 441--927. Nothing inside a retained excerpt is omitted or altered: every retained body is delivered exactly as the article's lines are written, so no omission receipt of any kind accompanies this package. Citations: the retained excerpts carry no citation command at all, so no bibliography entry of the article is transcribed and no reference is redistributed. Reference adaptation: none was needed and none was made; every reference inside the delivered unit resolves inside it, so no unresolved reference or citation is produced. Printed slips kept literal: no printed word, symbol, hypothesis or number of the counted statement or of its proof was altered. Layout and class: the article's own class is \texttt{amsart} with packages including babel, inputenc, amsfonts, amssymb, amsmath, caption, subcaption, siunitx, amsthm, graphicx, booktabs, xcolor, epsfig, fancyhdr; the wrapper is the lane's pinned \texttt{amsart} editorial wrapper at 10pt on A4, which restates the theorem-like environments the retained text needs and adds the article's own macro definitions that the retained text needs (none), with extra wrapper packages (bm, bbm, dsfont, xspace, cancel, mathtools). The delivered numbering is the unit's own. Assets: no retained span contains a figure environment, a graphics-inclusion command, a PDF-page inclusion, a table or a TikZ picture; no image file is copied into this package, the package declares no asset dependency and no image file is redistributed with it. Licence: version arXiv:2603.10096v1 of this article is distributed under the Creative Commons Attribution 4.0 International licence, as recorded in the arXiv licence metadata for this exact version and shown on the abstract page https://arxiv.org/abs/2603.10096v1. A full rights notice is delivered at licenses/arxiv-2603.10096-CC-BY-4.0.txt. Characters: every retained excerpt is pure ASCII, so the delivered engine, XeLaTeX with its default Latin Modern text font, reports no missing character.
\begin{equation}\label{ceq}
\dot H(t) = - \Pi_{\mathcal{S}_+}(\nabla_H E(H(t))) = -2 \alpha\, e(t)\, \Pi_{\mathcal{S}_+}(\nabla_H e(t)), \quad \alpha>0,
\end{equation}

\begin{lemma} \label{lemma 1}  
Let $H(t)\in \mathbb{R}^{2\times 2}$ be a differentiable real symmetric matrix valued function. Suppose that for some $t_{0}$ the smallest eigenvalue 
$\lambda_{\min}(t_{0})$ is simple. Let $v_{\min}(t)$ be the corresponding 
unit eigenvector chosen smoothly near $t_{0}$. Then the function 
$\lambda_{\min}(t)$ is differentiable at $t_{0}$ and satisfies
\begin{equation}\label{dotl}
\dot{\lambda}_{\min}(t_{0})
=
v_{\min}(t_{0})^{\top}\dot{H}(t_{0})\, v_{\min}(t_{0}).
\end{equation}
\end{lemma}

\begin{theorem}
Let $H(t)\in\mathbb{R}^{2\times 2}$ evolve according to \eqref{ceq}.  
If $H(0)\geq 0$, then
\[
H(t)\geq 0 \qquad \text{for all } t\ge 0.
\]
\end{theorem}

\begin{proof} Suppose $H(0) \geq 0.$ Then   $  \lambda_{min}(0) \geq 0.$ For $t>0 ,$ let $\lambda_{\min}(t)$ be the minimal eigenvalue of $H(t)$, with unit eigenvector
$v_{\min}(t)$ so that
\[
H(t)v_{\min}(t)=\lambda_{\min}(t)v_{\min}(t), 
\qquad \|v_{\min}(t)\|=1.
\] So, in order to prove $H(t)\geq 0$ for all $t$, it suffices to show  $\lambda_{\min}(t) \geq 0,$ for all $t.$\\
 By Lemma \eqref{lemma 1} we get,
\begin{equation}
\dot\lambda_{\min}(t)
= v_{\min}(t)^{\!\top}\, \dot H(t)\, v_{\min}(t).
\label{eq:lambdaderivative}
\end{equation}
Substituting from $\dot H(t)$ from the dynamics \eqref{ceq} yields
\begin{equation}
\dot\lambda_{\min}(t)
= -2\alpha\, e(t)\,
v_{\min}^{\!\top}(t)\, 
\Pi_{S^+}(\nabla_H e(t))\, 
v_{\min}(t).
\label{eq:lambdathm}
\end{equation}
Since $\Pi_{S^+}(\cdot)$ is the orthogonal projection onto the set of positive semidefinite matrices we have
\begin{equation}
v_{\min}^{\!\top}(t)\,
\Pi_{S^+}(\nabla_H e(t))\,
v_{\min}(t)
\ge 0.
\label{eq:nonneg}
\end{equation}
We now study the evolution of $\lambda_{\min}(t)$.\\

\paragraph{Case 1: $\lambda_{\min}(t) > 0$.}
Suppose that at some time $t_0$ we have $\lambda_{\min}(t_0) > 0$, so that $H(t_0)$ is strictly positive definite.  
From \eqref{eq:lambdathm}--\eqref{eq:nonneg}, the derivative $\dot{\lambda}_{\min}(t)$ is finite and may have either sign depending on the value of $e(t)$.Since $H(t)$ depends continuously on $t$, and the eigenvalues of a symmetric matrix depend continuously on its entries, the function $\lambda_{\min}(t)$ is continuous. Consequently, starting from a strictly positive value, $\lambda_{\min}(t)$ cannot cross into the negative region in finite time. Therefore, there exists an interval $[t_0,t_1)$ such that
\[
\lambda_{\min}(t) > 0 \quad \text{for all } t \in [t_0,t_1),
\]
and hence
\[
H(t) \in S^+ \quad \text{for all } t \in [t_0,t_1).
\]
 
\paragraph{Case 2: $\lambda_{\min}(t) = 0$.}
Suppose that at time $t$ we have $\lambda_{\min}(t)=0$, so that
$H(t)\in \partial S^+$. Let $v_{\min}(t)$ be a unit eigenvector
associated with $\lambda_{\min}(t)$, so that
\[
H(t)v_{\min}(t)=0.
\]

From \eqref{eq:lambdathm}, we have
\[
\dot{\lambda}_{\min}(t)
= -2\alpha e(t)\,
v_{\min}(t)^\top
\Pi_{\mathcal S^+}\big(\nabla_H e(t)\big)
v_{\min}(t).
\]
Define
\[
g(t)
:= v_{\min}(t)^\top
\Pi_{\mathcal S^+}\big(\nabla_H e(t)\big)
v_{\min}(t).
\]
Since $\Pi_{\mathcal S^+}(\cdot)$ is symmetric positive semidefinite,
it follows that
\[
g(t) \ge 0.
\] To establish invariance of $\mathcal S^+$, we examine the geometry of the cone
at the boundary point $H(t)$. The tangent cone to $\mathcal S^+$ at $H(t)$ is
given by
\[
T_{\mathcal S^+}(H(t))
=
\left\{
Z \in \mathbb{R}^{2\times 2}_{\mathrm{sym}}
\;\middle|\;
v^\top Z v \ge 0
\;\; \text{for all } v \in \ker H(t)
\right\}.
\]
Since $v_{\min}(t)\in \ker H(t)$ and
\[
v_{\min}(t)^\top \dot H(t)\, v_{\min}(t)
=
-2\alpha e(t)\, g(t),
\]
the only directions that could decrease $\lambda_{\min}(t)$ correspond to
negative values of $v_{\min}(t)^\top \dot H(t) v_{\min}(t)$.
However, the projection $\Pi_{\mathcal S^+}$ removes all components of
$\nabla_H e(t)$ that point outside $\mathcal S^+$, ensuring that
$\dot H(t)\in T_{\mathcal S^+}(H(t))$.
Therefore, no outward-pointing velocity is permitted at the boundary, and
$\lambda_{\min}(t)$ cannot become negative. Hence,
\[
H(t) \in \mathcal S^+ \quad \text{for all } t \ge 0.
\]


Case 3:  Suppose there exists a time $t_1 > 0$ such that $\lambda_{\min}(t_1) < 0$. By continuity of $\lambda_{\min}(t)$ and the initial condition $\lambda_{\min}(0) \ge 0$, there must exist a first time $t_0 \in [0,t_1)$ with $\lambda_{\min}(t_0) = 0$. At this time, $\dot \lambda_{\min}(t_0) = 0$, so the minimal eigenvalue cannot decrease further, contradicting the assumption that $\lambda_{\min}(t_1) < 0$. Hence, $\lambda_{\min}(t) < 0$ cannot occur.\\

Combining these cases, it shows that $
\lambda_{\min}(t) \ge 0 \quad \text{for all } t \ge 0, $ and therefore $H(t)$ remains positive semidefinite for all $t$.

\end{proof}

\end{document}
